\documentclass[10pt,a4paper]{amsart}
\usepackage[margin=19mm]{geometry}
\usepackage{amsmath,amssymb,mathtools}
\usepackage[scr=rsfs,cal=euler]{mathalfa}
\usepackage{xfrac}
\usepackage[unicode,hidelinks,bookmarks=false]{hyperref}
\newcommand{\ie}{i.e.\ }
\newcommand{\cf}{cf.\ }
\newcommand{\resp}{respectively\ }
\newcommand{\Subsection}{\subsection}
\providecommand{\nobreakdash}{\nobreak\hbox{-}}
\setlength{\emergencystretch}{2em}
\multlinegap0pt
\usepackage{bm}
\usepackage{bbm}
\usepackage{dsfont}
\usepackage{xspace}
\usepackage{cancel}
\usepackage{mathtools}
\usepackage{amsthm}
\makeatletter
\@ifundefined{definition}{\newtheorem{definition}{Definition}}{}
\@ifundefined{proposition}{\newtheorem{proposition}{Proposition}}{}
\makeatother
\title[Effects of semigroup properties on local embeddability]{Effects of semigroup properties on local embeddability}
\author{Dmitry Kudryavtsev}
\date{Source: arXiv:2510.21349v1 (CC BY 4.0)}
\begin{document}
\maketitle

\noindent\emph{Source and licence.} Effects of semigroup properties on local embeddability. Version arXiv:2510.21349v1 (CC BY 4.0), distributed under the Creative Commons Attribution 4.0 International licence, as recorded in the arXiv licence metadata for this exact version and shown on the abstract page https://arxiv.org/abs/2510.21349v1. Full rights notice: licenses/arxiv-2510.21349-CC-BY-4.0.txt.\par\smallskip

\noindent\textit{Editorial scope.} Counted unit: the Proposition whose source label is \texttt{None} in the article (source0.tex lines 368--370, source label \texttt{None}), delivered together with the article's own complete original proof of it (source0.tex lines 371--450, one \texttt{proof} environment of 80 lines). No mathematics is written, repaired or re-derived: statement and proof are the article's own text line for line. Required context retained uncounted: def\_lef (lines 78--82). Every definition, display and label of those lines is retained unchanged. Omitted from this package and not redistributed: 1--77, 83--367, 451--1372. Nothing inside a retained excerpt is omitted or altered: every retained body is delivered exactly as the article's lines are written, so no omission receipt of any kind accompanies this package. Citations: the retained excerpts carry no citation command at all, so no bibliography entry of the article is transcribed and no reference is redistributed. Reference adaptation: none was needed and none was made; every reference inside the delivered unit resolves inside it, so no unresolved reference or citation is produced. Printed slips kept literal: no printed word, symbol, hypothesis or number of the counted statement or of its proof was altered. Layout and class: the article's own class is \texttt{article} with packages including inputenc, babel, amsmath, amssymb, amsthm, enumerate, color, graphicx, lastpage209; the wrapper is the lane's pinned \texttt{amsart} editorial wrapper at 10pt on A4, which restates the theorem-like environments the retained text needs and adds the article's own macro definitions that the retained text needs (none), with extra wrapper packages (bm, bbm, dsfont, xspace, cancel, mathtools). The delivered numbering is the unit's own. Assets: no retained span contains a figure environment, a graphics-inclusion command, a PDF-page inclusion, a table or a TikZ picture; no image file is copied into this package, the package declares no asset dependency and no image file is redistributed with it. Licence: version arXiv:2510.21349v1 of this article is distributed under the Creative Commons Attribution 4.0 International licence, as recorded in the arXiv licence metadata for this exact version and shown on the abstract page https://arxiv.org/abs/2510.21349v1. A full rights notice is delivered at licenses/arxiv-2510.21349-CC-BY-4.0.txt. Characters: every retained excerpt is pure ASCII, so the delivered engine, XeLaTeX with its default Latin Modern text font, reports no missing character.
\begin{definition}\label{def_lef}
A (semi)group $S$ is called {\em locally embeddable into the class of finite (semi)groups} (an {\em LEF} (semi)group  for short) if for every finite subset $H$ of $S$ there exist a finite (semi)group $F_H$ and an injective function $f_H: H \rightarrow F_H$ such that for all  $x,y \in H$ with $xy \in H$ we have $(xy)f_H = (xf_H)(yf_H)$.

We call $(F_H, f_H)$ an {\em approximating pair} for $H$.
\end{definition}

\begin{proposition}
The semigroup $Q = Sg \langle a,b,c,e,x | xb = cx, ac = ca, ae=ea, ec=ce, xca = xe, aex = ax\rangle$ is LEF.
\end{proposition}

\begin{proof}
Denote by $H_n$ the subset of $Q$ consisting of all elements represented by words of length $n$ or less in the alphabet $\{a,b,c,e,x\}$, $n \ge 5$.

Let $w$ be any word in $\{a,b,c,e,x\}$. It has the form  $m_0 y_0 m_1 y_1 \cdots m_{k-1} y_{k-1} m_k$ where $k \ge 0$, $y_i$'s are the maximal subwords of the form $x$, $b^q$ or $x b^q$ for $q > 0$ and $m_i$'s are words in $\{a,c,e\}$. Denote $k$ by $s(w)$ and denote the generating relations in $Q$ by $\rho_0$.

Consider the following rewriting system.

\begin{enumerate}
\item $xb \rightarrow cx, ca \rightarrow ac, ea \rightarrow ae, ec \rightarrow ce$;

\item $x a^\alpha c^\beta \rightarrow x a^{\alpha - \beta} e^\beta$, $0< \beta \le \alpha$;
\item $x a^\alpha c^\beta \rightarrow x c^{\beta - \alpha} e^\alpha$, $0< \alpha < \beta$;

\item $ a^\alpha c^\beta e^\gamma x  \rightarrow a^\alpha c^\beta x $, $0<\alpha$, $0 \le \beta$, $0 < \gamma$;

\item $x e^\gamma x \rightarrow xex$, $1< \gamma$;
\item $x c^\beta e^\gamma x \rightarrow x c^{\beta} e x$, $0< \beta$, $1 <\gamma$;
\item $x a^\alpha c^\beta e^\gamma x \rightarrow x a^{\alpha - \beta} x$, $0\le \beta < \alpha$, $0\le \gamma$ with $\beta$ and $\gamma$ being not simultaneously zero;
\item $x a^\alpha c^\beta e^\gamma x\rightarrow x e x$, $0 <\alpha = \beta$, $0\le \gamma$;
\item $x a^\alpha c^\beta e^\gamma x\rightarrow x c^{\beta - \alpha} e x$, $0 <\alpha < \beta$, $0\le \gamma$;
\end{enumerate}

It is evident that this system is noetherian as every rule makes the rewritten word smaller in the lexicographical order induced by $a<c<e<b<x$. It is also locally confluent. To see this, consider a word $w$ such that we can apply two rules to its subwords, resulting in distinct words $w_1$ and $w_2$. Our goal is to demonstrate that there exists a word $w_0$ such that we can rewrite both $w_1$ and $w_2$ into $w_0$. If the rewritten subwords do not intersect, the statement is evident. Below we will consider $w = w' t w''$ where $t$ is the joint of the two intersecting subwords undergoing the rewritings and $w', w''$ are arbitrary. Correspondingly, $w_i = w' t_i w''$, $i \in \{0,1,2\}$. Note that for all rewritings where the intersection is the letter $x$ and neither of the rules is $xb \rightarrow cx$ the confluence is evident as these rules do not affect the $x$ and do not push letters to the either side of it. For the rest of the rewriting pairs see Appendix A.

Thus, we can establish a normal form for every element in $Q$.

Let $F_n$ be the semigroup defined as follows:

\begin{itemize}
\item It is generated by $a,b,c,d,e,x$ and a zero;
\item We have $xb = cx, ac = ca, ae=ea, ec=ce, xca = xe, aex = ax$ in $F_n$;
\item We also have $a = a^{2n+1}, b = b^{2n+1}, c=c^{2n+1}, e=e^{2n+1}$;
\item If  $s(w)>n$ in $w$, then $w$ represents the zero.
\end{itemize}

We claim that this is a finite semigroup and that the map $f_n$ which is a restriction of the natural homomorphism from $Q^{0}$ to $F_n$ onto $H_n$ satisfies the criteria of Definition \ref{def_lef}.

Consider an arbitrary word $w$. Recall that we can write it as  $m_0 y_0 m_1 y_1 \cdots m_{k-1} y_{k-1} m_k$. We claim that the rewritings corresponding to relations $xb = cx, ac = ca, ae=ea, ec=ce, xca = xe, aex = ax$ or $a = a^{2n+1}, b = b^{2n+1}, c=c^{2n+1}, e=e^{2n+1}$ (denote these $\rho$ for brevity) keep $s(w)$ constant. It follows from the fact that the existence of $x$'s is not affected by elements of $\rho$, while for $b$'s the only way to completely remove its power is to change $xb$ into $cx$, which retains the relative position of the corresponding $y_i$.

Additionally, due to commutativity and periodicity of $a,c,e$, there are at most $(2n+1)^3$ distinct forms of $m_i$. Similarly, there are only $2 \cdot (2n+1)$ distinct forms of $y_i$.

Furthermore, it is evident that for any two words $u,v$ we have $s(uv) \ge \max (s(u) , s(v))$ as concatenation does not decrease this value (only the bordering powers of $b$ can become a single part in the partition of $uv$).

This means that the words with $s(w) > n$ form an ideal in $\{a,b,c,e,x\}^*$, while all the other words correspond to a finite number of elements of $F_n$. Thus, $F_n$ is finite.

In order to see that $f_n$ separates the elements of $H_n$, consider the following rewriting system corresponding to the given relations of $F_n$.

\begin{enumerate}
\item $xb \rightarrow cx, ca \rightarrow ac, ea \rightarrow ae, ec \rightarrow ce$;


\item $a^{2n+1} \rightarrow a, b^{2n+1} \rightarrow b, c^{2n+1} \rightarrow c, e^{2n+1} \rightarrow e$;

\item $x a^\alpha c^\beta \rightarrow x a^{\alpha - \beta} e^\beta$, $0< \beta \le \alpha  \le 2n$;
\item $x a^\alpha c^\beta \rightarrow x c^{\beta - \alpha} e^\alpha$, $0< \alpha < \beta  \le 2n$;
\item $x a^\alpha c^\beta e^\gamma \rightarrow x c^{2n+\beta -\alpha} e^{\alpha+\gamma -2n}$,  $0< \alpha, \gamma \le 2n$, $0 \le \beta \le 2n$, $\alpha+\gamma > 2n, \gamma+\beta \le 2n$;
\item $x a^\alpha c^\beta e^\gamma \rightarrow x a^{2n+\alpha - \beta} e^{\beta+\gamma -2n}$,  $0< \beta, \gamma \le 2n$,  $0 \le \alpha \le 2n$, $\alpha+\gamma \le 2n, \gamma+\beta > 2n$;
\item $x a^\alpha c^\beta e^\gamma \rightarrow x a^{\alpha - \beta} e^{\beta+\gamma - 2n}$,  $0< \alpha,\beta, \gamma \le 2n$, $\alpha+\gamma, \gamma+\beta > 2n$, $\alpha > \beta$;
\item $x a^\alpha c^\beta e^\gamma \rightarrow x e^{\beta+\gamma - 2n}$,  $0< \alpha,\beta, \gamma \le 2n$, $\alpha+\gamma, \gamma+\beta > 2n$,  $\alpha = \beta$;
\item $x a^\alpha c^\beta e^\gamma \rightarrow x c^{\beta - \alpha} e^{\alpha+\gamma -2n}$,  $0< \alpha,\beta, \gamma \le 2n$, $\alpha+\gamma, \gamma+\beta > 2n$, $\beta > \alpha$;

\item $ a^\alpha c^\beta e^\gamma x  \rightarrow a^\alpha c^\beta x $, $0<\alpha, \gamma  \le 2n$, $0 \le \beta \le 2n$;

\item $x e^\gamma x \rightarrow xex$, $1< \gamma  \le 2n$;

\item $x c^\beta e^\gamma x \rightarrow x c^{\beta} e x$, $0< \beta < n$, $1 <\gamma \le 2n$;
\item $x c^\beta e^\gamma x \rightarrow x a^{2n-\beta} x$, $n\le \beta < 2n$, $0 <\gamma \le 2n$ and $x c^{2n} e^\gamma x \rightarrow x e x$, $0 <\gamma \le 2n$;

\item $x a^\alpha c^\beta e^\gamma x \rightarrow x a^{\alpha - \beta} x$, $0<\alpha - \beta \le n$, $0<  \alpha  \le 2n$, $0\le \beta,\gamma \le 2n$ with $\beta$ and $\gamma$ not simultaneously zero;
\item $x a^\alpha c^\beta e^\gamma x \rightarrow x c^{2n-\alpha+\beta} e x$, $n<\alpha - \beta < 2n$, $0< \alpha  \le 2n$, $0\le \beta, \gamma \le 2n$;

\item $x a^\alpha c^\beta e^\gamma x\rightarrow x e x$, $0 <\alpha = \beta \le 2n$, $0\le \gamma \le 2n$;
\item $x a^\alpha c^\beta e^\gamma x\rightarrow x c^{\beta - \alpha} e x$, $0 < \beta - \alpha < n$, $0< \beta, \alpha  \le 2n$, $0\le \gamma \le 2n$;
\item $x a^\alpha c^\beta e^\gamma x\rightarrow x a^{2n - \beta + \alpha} x$, $n \le \beta - \alpha < 2n$, $0< \beta, \alpha  \le 2n$, $0\le \gamma \le 2n$;
\end{enumerate}

It is evident that this system is noetherian as every rule makes the rewritten word smaller in the aforementioned lexicographical order. It is also locally confluent which can be demonstrated using an argument similar to the one for $Q$, with the supplementary table found in Appendix B.

Furthermore, it is clear that the normal forms of the elements of $H_n$ do not represent zero and they are also irreducible under this rewriting system. Thus, $f_n$ separates elements of $H_n$, which concludes the proof.
\end{proof}

\end{document}
